% ----- CHAUPAI 37 -----

\newpage

\subsection*{\deva{सप्तत्रिंशत्तम चौपाई} (Thirty-Seventh Chaupai)}
\addcontentsline{toc}{subsection}{\deva{सप्तत्रिंशत्तम चौपाई} (Thirty-Seventh Chaupai) :  \deva{जै जै...}}

\begin{minipage}[t]{0.55\textwidth}
\vspace{0pt}

{\large \linespread{1.5}\selectfont
\deva{जै जै जै हनुमान गोसाईं।}\\
\deva{कृपा करहु गुरु देव की नाईं॥}
\par}

\vspace{0.5em}

{\large \linespread{1.5}\selectfont
\telu{జై జై జై హనుమాన గోసాయీఁ।}\\
\telu{కృపా కరహు గురు దేవ కీ నాయీఁ॥}
\par}

\vspace{0.5em}

{\large \linespread{1.3}\selectfont
\textit{jai jai jai hanumāna gosāī̃ |}\\
\textit{kṛpā karahu guru deva kī nāī̃ ||}
\par}
\end{minipage}%
\hfill
\begin{minipage}[t]{0.42\textwidth}
\vspace{0pt}
\begin{tcolorbox}[anchorbox]
\textbf{The Ultimate Mentor:} By accepting an ideal role model as our guide (\textit{Gurudeva}), we are asking for the unvarnished truth and the strict discipline required to elevate our lives to their highest potential. A true mentor challenges us to grow.
\vspace{0.5em}\newline By accepting an ideal role model, we invite the strict discipline needed to universally elevate our own potential.\newline\null\hfill\href{https://www.valmiki.iitk.ac.in/sloka?field_kanda_tid=5&language=dv&field_sarga_value=1}{\textit{Sundara Kanda, 1:205}}
\end{tcolorbox}
\end{minipage}

\vspace{0.5em}
\subsubsection*{\textbf{Visual Rhythm Map:}}
\begin{table}[h!]
\centering
\renewcommand{\arraystretch}{1.2}
\noindent\makebox[\textwidth][l]{%
  {%
  \setlength{\tabcolsep}{0pt}%
  \begin{tabularx}{1.0625\textwidth}{|*{17}{>{\centering\arraybackslash\hsize=1.0\hsize}X|}}
  \hline
  \multicolumn{2}{|>{\centering\arraybackslash\hsize=2.0\hsize}X|}{\textbf{\guru{जै}} \par (2)} &
  \multicolumn{2}{>{\centering\arraybackslash\hsize=2.0\hsize}X|}{\textbf{\guru{जै}} \par (2)} &
  \multicolumn{2}{>{\centering\arraybackslash\hsize=2.0\hsize}X|}{\textbf{\guru{जै}} \par (2)} &
  \multicolumn{5}{>{\centering\arraybackslash\hsize=5.0\hsize}X|}{\textbf{\laghu{ह}\laghu{नु}\guru{मा}\laghu{न}} \par (5)} &
  \multicolumn{6}{>{\centering\arraybackslash\hsize=6.0\hsize}X|}{\textbf{\guru{गो}\guru{सा}\guru{ईं}} \par (6)} \\
  \hline
  \end{tabularx}%
  }%
  \hspace{0.01\textwidth}%
  \begin{minipage}{0.1\textwidth}
  \vspace{-0.5em}
  \centering
  {\Huge \textbf{17}}
  \end{minipage}%
}
\vspace{-1.5em}
\end{table}

\begin{table}[h!]
\centering
\renewcommand{\arraystretch}{1.2}
\noindent\makebox[\textwidth][l]{%
  {%
  \setlength{\tabcolsep}{0pt}%
  \begin{tabularx}{1.0625\textwidth}{|*{17}{>{\centering\arraybackslash\hsize=1.0\hsize}X|}}
  \hline
  \multicolumn{3}{|>{\centering\arraybackslash\hsize=3.0\hsize}X|}{\textbf{\laghu{कृ}\guru{पा}} \par (3)} &
  \multicolumn{3}{>{\centering\arraybackslash\hsize=3.0\hsize}X|}{\textbf{\laghu{क}\laghu{र}\laghu{हु}} \par (3)} &
  \multicolumn{2}{>{\centering\arraybackslash\hsize=2.0\hsize}X|}{\textbf{\laghu{गु}\laghu{रु}} \par (2)} &
  \multicolumn{3}{>{\centering\arraybackslash\hsize=3.0\hsize}X|}{\textbf{\guru{दे}\laghu{व}} \par (3)} &
  \multicolumn{2}{>{\centering\arraybackslash\hsize=2.0\hsize}X|}{\textbf{\guru{की}} \par (2)} &
  \multicolumn{4}{>{\centering\arraybackslash\hsize=4.0\hsize}X|}{\textbf{\guru{ना}\guru{ईं}} \par (4)} \\
  \hline
  \end{tabularx}%
  }%
  \hspace{0.01\textwidth}%
  \begin{minipage}{0.1\textwidth}
  \vspace{-0.5em}
  \centering
  {\Huge \textbf{17}}
  \end{minipage}%
}
\vspace{-1.5em}
\end{table}


\vspace{1em}
\textbf{ \deva{सरल-संस्कृतम्}:}\\
 \deva{हे गोस्वामिन् (इन्द्रियाणां स्वामिन्) हनुमन्! भवतः जयः, जयः, जयः अस्तु।}\\
 \deva{भवान् मम उपरि इष्ट-गुरुदेवः इव कृपां करोतु॥}\\


Hail, Hail, All Hail to you, Lord Hanuman, Master of the Senses!\\
Please bestow your grace upon me, just like a Supreme Guru.


\vspace{1em}
\newpage
\textbf{\deva{प्रतिपदार्थः} (Meaning of each word):}
\begin{table}[h!]
\centering
\renewcommand{\arraystretch}{1.2}
\begin{tabularx}{\textwidth}{@{} l l X X @{}}
\toprule
\textbf{Awadhi} & \textbf{Sanskrit} & \textbf{English} & \textbf{Grammar \& Notes} \\
\midrule
\deva{जै जै जै} / \telu{జై జై జై}  &  \deva{जय, जय, जय}  &  Victory / Hail (x3)  &   \\
\deva{हनुमान} / \telu{హనుమాన}  &  \deva{हनुमन्}  &  Hanuman  &   \\
\deva{गोसाईं} / \telu{గోసాయీఁ}  &  \deva{गोस्वामिन्}  &  Lord of Senses  & Go (Senses) + Swami (Lord) \\
\deva{कृपा करहु} / \telu{కృపా కరహు}  &  \deva{कृपां करोतु}  &  Have mercy / grace  & Imperative `-hu` suffix \\
\deva{गुरु देव} / \telu{గురు దేవ}  &  \deva{गुरूदेवः}  &  Supreme Guru  &   \\
\deva{की नाईं} / \telu{కీ నాయీఁ}  &  \deva{इव / सदृशम्}  &  Like / In the manner of  & Pure Awadhi idiom \\
\bottomrule
\end{tabularx}
\end{table}

\begin{tcolorbox}[poetrybox]
\textbf{The 17-Beat Anomaly \& Bhava-Pradhan:}\\
Both lines mathematically sum up to 17 beats! However, because this is a \textit{Bhava Pradhan} poem, the strict 16-beat limit gracefully bows to devotion. The joyful \deva{जय} naturally compresses into an unbroken triplet (\textit{je-je-jai}) in Line 1, and the short 'u' in \deva{करहु} or \deva{गुरु} slips by silently in Line 2 without interrupting the singing flow. These organic inconsistencies are beautiful mnemonic markers that enhance the chanting experience!
\end{tcolorbox}

\newpage
